\section*{Abkürzungsverzeichnis und Begriffsdefinitionen}
\addcontentsline{toc}{section}{Abkürzungsverzeichnis und Begriffsdefinitionen}

\subsection*{Abkürzungsverzeichnis}

\renewcommand{\arraystretch}{1.2}
\begin{xltabular}{\textwidth}{p{3.0cm}X}
\textbf{Abkürzung} & \textbf{Bedeutung} \\
\hline

API   & Application Programming Interface \\
CLI   & Command Line Interface \\
CSV   & Comma-Separated Values \\
IEM   & Internal Engineering Model \\
KI    & Künstliche Intelligenz \\
LLM   & Large Language Model \\
MBSE  & Model-Based Systems Engineering \\
OCR   & Optical Character Recognition \\
PLM   & Product Lifecycle Management \\
QVT   & Query/View/Transformation \\
REDPA & Reusable Engineering Document Processing Architecture \\
RFLP  & Requirements, Functional, Logical, Physical \\
SOP   & Standard Operating Procedure \\
SysML & Systems Modeling Language \\
W3C   & World Wide Web Consortium \\
XMI   & XML Metadata Interchange \\

\end{xltabular}

\vspace{0.8cm}

\subsection*{Begriffsdefinitionen}

\renewcommand{\arraystretch}{1.25}
\begin{xltabular}{\textwidth}{p{4.2cm}X}
\textbf{Begriff} & \textbf{Verwendung in dieser Arbeit} \\
\hline

Brownfield-Transformation &
Überführung bereits vorhandener Engineering-Information in einen
modellbasierten Kontext unter Berücksichtigung bestehender Systeme,
Dokumente, Datenbestände und Informationsstrukturen. \\

Engineering-Quelle &
Rolle eines Dokuments, Datensatzes oder sonstigen Informationsbestands als
fachliche Ausgangsbasis des Transformationsprozesses. Eine
Engineering-Quelle darf quellengebundene Evidenz über das betrachtete System
begründen. \\

Kontextquelle &
Rolle einer Informationsquelle, die zusätzliche Informationen zur
Interpretation oder Einordnung bereitstellt. Eine Kontextquelle unterstützt
die Verarbeitung, begründet selbst jedoch keine fachliche Evidenz über das
betrachtete System. \\

Quellengebundene Evidenz &
Aus einer Engineering-Quelle identifizierte Information, deren konkrete
Herkunft erhalten bleibt und die als fachliche Grundlage nachfolgender
Interpretations- und Autorisierungsentscheidungen dient. \\

Engineering Subject &
Quellengebundener fachlicher Referenzgegenstand, auf den unterschiedliche
Persona-Perspektiven angewendet und deren Interpretationsergebnisse
miteinander verglichen werden können. \\

Persona &
Konfigurierte fachliche Perspektive, aus der ein bestimmter
Interpretations- oder Modellierungsschritt durchgeführt wird. Mehrere
Personas können denselben fachlichen Gegenstand mit unterschiedlichen
Schwerpunkten betrachten. \\

Human Authority &
Oberbegriff für die explizite menschliche Entscheidungsbefugnis an einem
definierten fachlichen oder modellbezogenen Übergang des
Transformationsprozesses. \\

Human Engineering Authority &
Menschliche Entscheidungsbefugnis zur Autorisierung der fachlichen Bedeutung
und weiteren Verwendung von Engineering-Information. \\

Engineering Meaning &
Fachliche Bedeutung einer autorisierten Engineering-Information unabhängig
von ihrer konkreten strukturellen Repräsentation im Zielmodell. \\

Model Candidate &
Vorschlag für ein mögliches Modellelement, eine Beziehung oder eine
strukturelle Zuordnung, der aus autorisierter Engineering-Information
abgeleitet und als Entscheidungsgrundlage für die Model Authority
bereitgestellt wird. \\

Model Authority &
Menschliche Entscheidungsbefugnis zur Autorisierung der strukturellen
Einordnung und Repräsentation einer Engineering-Information im Zielmodell. \\

Internal Engineering Model (IEM) &
Interne Modellrepräsentation, in der autorisierte Modellinhalte unter Erhalt
ihrer fachlichen Bedeutung und strukturellen Einordnung zusammengeführt
werden. Das IEM bildet die Grundlage für die nachgelagerte
zielmodellspezifische Formulierung und SysML-v2-Erzeugung. \\

Projektzulassung (\emph{Project Fit}) &
Kontrollierter Bewertungsschritt, der prüft, ob eine Engineering-Quelle
plausibel dem betrachteten System- und Projektkontext zugeordnet und für die
projektweite Verarbeitung berücksichtigt werden kann. \\

Reusable Engineering Document Processing Architecture (REDPA) &
Im Ausblick verwendeter Arbeitsbegriff für einen möglichen
wiederverwendbaren Kern der entwickelten Informationsverarbeitungsarchitektur.
REDPA überträgt insbesondere generische Mechanismen der Quellenverarbeitung,
Provenienz, persistenten Verarbeitungszustände und menschlichen
Autorisierung auf weitere dokumentenzentrierte Engineering-Aufgaben, während
fach- und zielspezifische Bestandteile angepasst oder ersetzt werden können. \\

Akzeptanzniveau &
Im Ausblick vorgeschlagene Bewertungsgröße für lernende Erweiterungen des
Transformationssystems. Sie beschreibt, in welchem Umfang maschinell erzeugte
Vorschläge durch die zuständige Human Authority ohne oder mit nur geringer
fachlicher Änderung übernommen werden. \\

Provenienz &
Nachvollziehbare Herkunft und Entstehung einer Information einschließlich
der relevanten Verarbeitungs- und Entscheidungszustände, aus denen sie
hervorgegangen ist. \\

Nachverfolgbarkeit (\emph{Traceability}) &
Nachvollziehbare Verknüpfung zwischen Engineering-Quellen,
Verarbeitungszuständen, Autorisierungsentscheidungen, Modellzuständen und
resultierenden Modellartefakten. \\

Persistenz &
Dauerhafte Speicherung relevanter Informations-, Entscheidungs- und
Modellzustände. Frühere Zustände bleiben dabei erhalten und können mit
nachfolgenden Zuständen und Entscheidungen verknüpft werden. \\

Validierung &
Prüfung eines erzeugten SysML-v2-Artefakts anhand festgelegter formaler und
technischer Kriterien sowie seiner Verarbeitbarkeit in der verwendeten
SysML-v2-Toolumgebung. \\

Fail-closed &
Verarbeitungsprinzip, bei dem ein Verarbeitungsschritt nur ausgeführt wird,
wenn die dafür erforderlichen Voraussetzungen erfüllt sind. Fehlen
notwendige Autorisierungen, Referenzen oder eindeutig anwendbare
Transformationsregeln, wird die weitere Verarbeitung an dieser Stelle
blockiert. \\

\end{xltabular}